% AY2025 材料力学 〔II〕（転記）
% 出典: 04_材料力学/original/04_材料力学/2025/「2025za - コピー (2).pdf」
% 要約: 直角に折れ曲がった円形断面のはり(直径d)。点A固定, 区間AB(長さL) → 直角に折れ →
%       区間BC(長さL/2) → 点Cに下向き集中荷重P。縦弾性係数E, 横弾性係数G。
%       点Dは AからL/4の断面, u は断面D上の点。
%       (1) 固定端A支持力RA・支持モーメントMA・支持トルクTA, (2) 区間ABのSFD/BMD,
%       (3) AからL/4の断面Dの点uのσx,σy,τxy, (4) 点BたわみδB・ねじれ角ψB・点CたわみδC
%       (たわみ基礎式, 断面二次I, 断面二次極Ip)。
% rem: 原本は3次元俯瞰図。丸棒(x方向, 区間AB)+丸棒(y方向, 区間BC)のL字構成を簡略立体図で再現する。

\section*{〔II〕}

図2のように, 直角に折れ曲がった円形断面のはり（直径 $d$）の点 A が固定され, 点 C に下向きの
集中荷重 $P$ が作用する. 縦弾性係数を $E$, 横弾性係数を $G$ とし, 次の問い (1)\,$\sim$\,(4) に答えよ.

\begin{center}
\begin{tikzpicture}[>=Latex]
  % 固定壁（A, z-up 斜め平面）
  \fill[pattern=north east lines] (0.2,0.6) -- (0.2,3.0) -- (1.2,3.4) -- (1.2,1.0) -- cycle;
  \draw (0.2,0.6) -- (0.2,3.0) -- (1.2,3.4) -- (1.2,1.0) -- cycle;
  % はり 区間AB（x方向, 下右へ, 長さL）
  \draw[line width=11pt,gray!25,line cap=round] (0.9,2.4) -- (6.0,1.0);
  % 区間BC（y方向, 右へ, 長さL/2）
  \draw[line width=11pt,gray!25,line cap=round] (6.0,1.0) -- (8.3,1.6);
  \draw[fill=gray!25] (8.3,1.6) ellipse (0.30 and 0.11);
  % 軸（A点）
  \draw[->] (0.9,2.4) -- (0.9,3.7) node[left]{$z$};
  \draw[->] (0.9,2.4) -- (2.0,2.95) node[above]{$y$};
  \draw[->] (6.0,1.0) -- (7.0,1.0) node[below]{$x$};
  % 節点
  \node[left] at (0.8,2.45) {A}; \node[below] at (6.0,0.7) {B}; \node[right] at (8.4,1.65) {C};
  % 点 u, 断面D（L/4）
  \fill (2.3,2.25) circle (1.8pt); \node[above] at (2.3,2.3) {u};
  \draw[dashed] (2.3,1.95) ellipse (0.1 and 0.3);
  \node[below] at (2.3,1.75) {D};
  % 荷重 P（C, 下向き）
  \draw[->,very thick] (8.3,2.7) -- (8.3,1.75); \node[right] at (8.35,2.5) {$P$};
  % 寸法 L（A-B）, L/4（A-D）, L/2（B-C）
  \draw[<->] (1.1,3.05) -- (6.0,1.65) node[midway,above,font=\scriptsize]{$L$};
  \draw[<->] (1.0,1.95) -- (2.3,1.6) node[midway,below,font=\scriptsize]{$L/4$};
  \draw[<->] (6.1,0.65) -- (8.3,1.25) node[midway,below,font=\scriptsize]{$L/2$};
  % 断面D インセット（左下）
  \begin{scope}[shift={(0.9,-1.6)}]
    \draw (0,0) circle (0.55);
    \draw[->] (0,-0.85) -- (0,0.95) node[right]{$z$};
    \draw[->] (-0.85,0) -- (0.95,0) node[above]{$y$};
    \fill (0,0.55) circle (1.8pt); \node[above left] at (0,0.55) {u};
    \node at (0,-1.25) {\underline{断面 D}};
  \end{scope}
\end{tikzpicture}

\vspace{1mm}
図2
\end{center}

\begin{enumerate}[label=(\arabic*)]
  \item 固定端 A での支持力 $R_A$, 支持モーメント $M_A$, 支持トルク $T_A$ を図示し, 数式で示せ.

  \item 区間 AB でのせん断力図（SFD）と曲げモーメント図（BMD）を求めて図示せよ.

  \item 固定端 A から $L/4$ の位置にある断面 D の点 u の $x$ 方向の応力 $\sigma_x$,
        $y$ 方向の応力 $\sigma_y$, せん断応力 $\tau_{xy}$ を求めよ.

  \item たわみの基礎式を用いて, 点 B のたわみ $\delta_B$, ねじれ角 $\psi_B$, および点 C の
        たわみ $\delta_C$ を求めよ. ただし, はりの断面二次モーメントは $I$,
        断面二次極モーメントは $I_p$ として解答せよ.
\end{enumerate}
